\documentclass[11pt,a4paper]{article}

\usepackage[margin=25mm]{geometry}
\usepackage{amsmath,amssymb,bm,booktabs,array,longtable,tabularx}
\usepackage{graphicx}
\usepackage{hyperref}
\usepackage{siunitx}
\usepackage{xcolor}

\newcommand{\Tens}{\bm P}
\newcommand{\vect}[1]{\bm{#1}}
\newcommand{\Tr}{\operatorname{Tr}}
\newcommand{\rank}{\operatorname{rank}}
\newcommand{\diag}{\operatorname{diag}}
\newcommand{\cond}{\operatorname{cond}}
\newcommand{\pseudoDet}{\operatorname{pdet}}
\newcommand{\Null}{\mathcal N}
\newcommand{\Range}{\mathcal R}
\newcolumntype{Y}{>{\raggedright\arraybackslash}X}

\newcommand{\generatedfigure}[3]{%
  \begin{figure}[htbp]
    \centering
    \IfFileExists{figures/#1}{%
      \includegraphics[width=0.88\linewidth]{figures/#1}%
    }{%
      \fbox{\parbox[c][45mm][c]{0.82\linewidth}{\centering
      \textcolor{red!70!black}{\textsf{[generated input pending: }\nolinkurl{figures/#1}\textsf{]}}}}%
    }
    \caption{#2}
    \label{#3}
  \end{figure}%
}

\title{Identifiability of Deuteron Tensor Polarization\\
with One and Two L/R/U/D Polarimeters}
\author{DPOLAR analysis note}
\date{Generated from repository conventions; numerical tables are reproducible analysis products}

\begin{document}
\maketitle

\begin{abstract}
This note separates three questions that are often conflated in discussions of
deuteron polarimetry: whether a tensor component appears in the cross section,
whether the detector geometry provides an independent response to it, and
whether that response survives profiling of luminosity and detector nuisance
parameters.  The analysis uses the Cartesian spin--1 convention already used
by the repository.  An ideal single L/R/U/D ring has one exact tensor-geometry
null direction, $p_{xy}$, when absolute normalization is known.  With unknown
global normalization its common-mode $p_{zz}$ response is also lost at one
polar angle.  Multiple polar angles can recover $p_{zz}$ only when their
acceptance-averaged $A_{zz}$ responses differ and share a constrained
normalization.  A second station increases structural rank only when known spin
transport rotates one station's null space away from the other's.  Around an
axial vertical $p_{yy}$ state, one infinitesimal tilt is visible at first order
and the orthogonal tilt is invisible in a perfectly cardinal four-arm geometry.
These results support an after-SRC polarimeter as a direct monitor of the tensor
polarization delivered by the accelerator chain, while making the scientific
case for a second station conditional on a specific spin-transport or
contamination measurement.
\end{abstract}

\section{Scope, authority, and limitations}

The source of truth for polarization normalization is the repository's
Cartesian density-matrix convention.  The source of truth for the measured
reaction response is the repository's polarized $d$--$p$ elastic-scattering
formula.  The current tabulation contains $iT_{11}$, $T_{20}$, $T_{21}$, and
$T_{22}$ as functions of deuteron center-of-mass scattering angle.  The legacy
C++ scalar count path presently consumes only $T_{20}$ and $T_{22}$; the full
response requires all four observables.

All rank claims below are analytic unless explicitly labeled as generated
results.  Scenario-specific singular values, Fisher eigenvalues, uncertainties,
and figures must be inserted from the machine-readable output of the analysis
executable.  Placeholders in this note are deliberate and must not be read as
physics forecasts.

The current repository does not contain a complete, sourceable, element-by-
element spin-transfer map from the proposed after-SRC station to the proposed
pre-SAMURAI station.  Rotation scans are therefore generic sensitivity studies,
not predictions of the RIBF beam line.

\section{Spin--1 tensor degrees of freedom}

In the ordered $S_z$ basis, define the Hermitian Cartesian spin--1 tensor
operators
\begin{equation}
 Q_{ij}=\frac{3}{2}(S_iS_j+S_jS_i)-2\delta_{ij}I,
 \qquad
 p_i=\Tr(\rho S_i),
 \qquad
 p_{ij}=\Tr(\rho Q_{ij}).
\end{equation}
The density matrix is expanded as
\begin{equation}
 \rho=\frac{I}{3}
 +\frac12\sum_i p_iS_i
 +\frac19\sum_i p_{ii}Q_{ii}
 +\frac29\sum_{i<j}p_{ij}Q_{ij}.
 \label{eq:rho-expansion}
\end{equation}
The tensor polarization is the real symmetric traceless matrix
\begin{equation}
 \Tens=
 \begin{pmatrix}
 p_{xx}&p_{xy}&p_{xz}\\
 p_{xy}&p_{yy}&p_{yz}\\
 p_{xz}&p_{yz}&p_{zz}
 \end{pmatrix},
 \qquad
 p_{xx}+p_{yy}+p_{zz}=0.
 \label{eq:tensor}
\end{equation}
It therefore has five independent degrees of freedom.  For a state diagonal in
magnetic substates about an axis $a$,
\begin{equation}
 p_a=N_{+1}-N_{-1},
 \qquad
 p_{aa}=N_{+1}+N_{-1}-2N_0=1-3N_0.
\end{equation}
For zero vector polarization the axial tensor range is
$-2\leq p_{aa}\leq1$.  The legacy scalar scan range $0\leq p\leq1$ is a
scenario choice and must not be mistaken for the physical tensor range.

\subsection{Orthonormal numerical basis and physical reporting basis}

A well-conditioned internal basis is orthonormal under the Frobenius inner
product, $\langle A,B\rangle=\Tr(A^{\mathsf T}B)$:
\begin{align}
 E_1&=\frac{1}{\sqrt2}\diag(1,-1,0),&
 E_2&=\frac{1}{\sqrt2}
 \begin{pmatrix}0&1&0\\1&0&0\\0&0&0\end{pmatrix},&
 E_3&=\frac{1}{\sqrt2}
 \begin{pmatrix}0&0&1\\0&0&0\\1&0&0\end{pmatrix},\\
 E_4&=\frac{1}{\sqrt2}
 \begin{pmatrix}0&0&0\\0&0&1\\0&1&0\end{pmatrix},&
 E_5&=\frac{1}{\sqrt6}\diag(-1,-1,2).&&
\end{align}
Writing $q_a=\Tr(\Tens E_a)$ gives
\begin{equation}
 \vect q=
 \begin{pmatrix}
 (p_{xx}-p_{yy})/\sqrt2\\
 \sqrt2\,p_{xy}\\
 \sqrt2\,p_{xz}\\
 \sqrt2\,p_{yz}\\
 \sqrt{3/2}\,p_{zz}
 \end{pmatrix}.
 \label{eq:orthonormal-coordinates}
\end{equation}
The numerical basis must not replace the experiment-facing reporting basis
\begin{equation}
 \vect r=(p_{xx}-p_{yy},\ p_{xy},\ p_{xz},\ p_{yz},\ p_{zz})^{\mathsf T}.
 \label{eq:physical-basis}
\end{equation}
The inverse relations are
\begin{equation}
 p_{xx}-p_{yy}=\sqrt2q_1,\qquad
 (p_{xy},p_{xz},p_{yz})=\frac{1}{\sqrt2}(q_2,q_3,q_4),\qquad
 p_{zz}=\sqrt{\frac23}q_5,
\end{equation}
with $p_{xx}=[(p_{xx}-p_{yy})-p_{zz}]/2$ and
$p_{yy}=[-(p_{xx}-p_{yy})-p_{zz}]/2$.

\section{Polarized elastic-scattering response}

In the local incident-beam frame used by the repository,
\begin{align}
 \frac{\sigma(\theta,\phi)}{\sigma_0(\theta)}=1
 &+\frac32(p_x\sin\phi+p_y\cos\phi)A_y(\theta)\notag\\
 &+\frac23(p_{xz}\cos\phi-p_{yz}\sin\phi)A_{xz}(\theta)\notag\\
 &+\frac16\left[(p_{xx}-p_{yy})\cos2\phi
 -2p_{xy}\sin2\phi\right]
 [A_{xx}(\theta)-A_{yy}(\theta)]\notag\\
 &+\frac12p_{zz}A_{zz}(\theta).
 \label{eq:cross-section}
\end{align}
The repository's spherical-to-Cartesian conversion is
\begin{equation}
 A_y=\frac{2}{\sqrt3}iT_{11},\qquad
 A_{xz}=-\sqrt3T_{21},\qquad
 A_{xx}-A_{yy}=2\sqrt3T_{22},\qquad
 A_{zz}=\sqrt2T_{20}.
 \label{eq:analyzing-conversion}
\end{equation}
Equations~\eqref{eq:cross-section} and \eqref{eq:analyzing-conversion}
reproduce the maintained scalar weights
\begin{align}
 f_{p_{zz}}&=1+\frac{\sqrt2}{2}T_{20}p_{zz},\\
 f_{LR/UD}&=1\mp\frac{\sqrt3}{2}T_{22}p_{yy}
 -\frac{\sqrt2}{4}T_{20}p_{yy}
\end{align}
for a longitudinal axial tensor and a vertical axial tensor, respectively.

An older displayed closed-form LR/UD inversion contains a factor-of-two error
in its $A_{zz}$ denominator.  From Eq.~\eqref{eq:cross-section} and the
maintained sector weights, the cardinal-ring result is
\begin{equation}
 R_{LRUD}=
 \frac{N_L+N_R-N_U-N_D}{N_L+N_R+N_U+N_D}
 =-\frac{p_{yy}(A_{xx}-A_{yy})}{4-p_{yy}A_{zz}}.
 \label{eq:lrud-corrected}
\end{equation}
The forward-model likelihood follows the maintained sector weights and is not
replaced by the obsolete analytic denominator.

The signs of the $U/D$ first harmonics in an older expanded set of paper
equations disagree with Eq.~\eqref{eq:cross-section}.  This analysis adopts the
general equation and the maintained response implementation: $L,U,R,D$ are at
$\phi=(0,\pi/2,\pi,3\pi/2)$, so $U-D$ has a negative $p_{yz}A_{xz}$ sign.
Changing the experimental azimuth convention requires changing both channel
metadata and interpretation, not only plot labels.

\section{Acceptance-integrated channel means}

For detector channel $k$, the expected Poisson mean is
\begin{equation}
 \mu_k(\vect q,\vect\eta)=
 \mathcal L\,\epsilon_k
 \int_{\Omega_k}d\Omega\,
 \sigma_0(\theta)
 \left[1+\sum_{a=1}^{5}c_a(\theta,\phi)q_a
 +\sum_{v=1}^{3}d_v(\theta,\phi)p_v\right],
 \label{eq:channel-mean}
\end{equation}
where $d\Omega=\sin\theta\,d\theta\,d\phi$, $\mathcal L$ contains beam dose
and target areal density, $\epsilon_k$ is a channel efficiency, and the
coefficient functions are obtained from Eq.~\eqref{eq:cross-section} and the
basis conversion.  Define the unpolarized integral and tensor response
coefficients
\begin{align}
 b_k&=\epsilon_k\int_{\Omega_k}d\Omega\,\sigma_0(\theta),\\
 H_{ka}&=\epsilon_k\int_{\Omega_k}d\Omega\,
 \sigma_0(\theta)c_a(\theta,\phi).
\end{align}
Then tensor-only means are linear,
\begin{equation}
 \vect\mu=\mathcal L(\vect b+H\vect q).
 \label{eq:linear-means}
\end{equation}
The physical target is an integral over the actual detector face.  Center-angle
evaluation is used only for the analytic ring check.  The numerical study in
this note instead follows the maintained finite rectangular top-hat acceptance,
including analytic integration of the azimuthal harmonics across each opening;
the approximation is stated in every output.

\subsection{Legacy acceptance caution}

The maintained scalar C++ count integrator presently multiplies a midpoint
sum in $\theta_{\rm cm}$ by a constant $\Delta\phi$ but omits the
$\sin\theta_{\rm cm}$ solid-angle factor.  A separate coverage-map path does
include $\sin\theta\,d\theta\,d\phi$, but that map is not used by scalar count
inference.  Consequently:
\begin{enumerate}
 \item existing scalar numerical outputs remain regression anchors and must not
 be changed silently;
 \item the present general response deliberately preserves the legacy
 $d\theta\,d\phi$ measure while adding finite-azimuth harmonics, so its absolute
 rate forecasts are provisional;
 \item absolute rates from the legacy and corrected measures must not be mixed;
 \item a future exact circular-face integration with $\sin\theta$ must be a
 separately named acceptance mode and reported with before-and-after regression
 tests.
\end{enumerate}

\section{Response matrix and local identifiability}

For parameters $\theta_a$, define
\begin{equation}
 J_{ka}=\frac{\partial\mu_k}{\partial\theta_a}.
\end{equation}
For tensor coordinates, $J_{ka}=\mathcal L H_{ka}$.  If the luminosity nuisance
is parameterized as $\lambda=\log\mathcal L$, then
$\partial\mu_k/\partial\lambda=\mu_k$.  A log parameter avoids scale-dependent
finite differences and enforces positive luminosity.

For the response in the orthonormal basis of Eq.~\eqref{eq:orthonormal-coordinates},
with no additional column scaling, compute
\begin{equation}
 J=U\Sigma V^{\mathsf T}.
\end{equation}
The reported output must contain all singular values and right-singular vectors,
not only an integer rank.  A singular value is counted when
\begin{equation}
 \sigma_i>\tau,\qquad
 \tau=\max(\tau_{\rm abs},\tau_{\rm rel}\sigma_1).
 \label{eq:rank-threshold}
\end{equation}
A suitable default after documented column scaling is
$\tau_{\rm rel}=10^{-10}$ and $\tau_{\rm abs}=10^{-12}$, accompanied by a
threshold scan.  Rank conclusions must be stable over at least two orders of
magnitude around the chosen relative threshold.  Every null or near-null right
singular vector must be converted from $\vect q$ to the physical basis
\eqref{eq:physical-basis} and normalized with a stated sign convention.

\section{Analytic sanity check: one ideal L/R/U/D ring}

Let $B=A_{xx}-A_{yy}$ and use the physical tensor coordinates
\eqref{eq:physical-basis}.  At a fixed polar angle, the polarization part of
the response row at azimuth $\phi$ is
\begin{equation}
 h(\phi)=
 \left(
 \frac{B}{6}\cos2\phi,
 -\frac{B}{3}\sin2\phi,
 \frac{2A_{xz}}{3}\cos\phi,
 -\frac{2A_{xz}}{3}\sin\phi,
 \frac{A_{zz}}{2}
 \right).
 \label{eq:ideal-row}
\end{equation}
For $L,U,R,D$ this becomes
\begin{equation}
 H_{\rm ring}=C
 \begin{pmatrix}
 B/6&0&2A_{xz}/3&0&A_{zz}/2\\
 -B/6&0&0&-2A_{xz}/3&A_{zz}/2\\
 B/6&0&-2A_{xz}/3&0&A_{zz}/2\\
 -B/6&0&0&2A_{xz}/3&A_{zz}/2
 \end{pmatrix},
 \label{eq:ring-matrix}
\end{equation}
where $C$ is the common unpolarized acceptance integral.

\subsection{Geometry null and measured harmonics}

Equation~\eqref{eq:ring-matrix} gives, without hard-coded null vectors,
\begin{align}
 L-R&\propto\frac43p_{xz}A_{xz},\\
 U-D&\propto-\frac43p_{yz}A_{xz},\\
 (L+R)-(U+D)&\propto\frac23(p_{xx}-p_{yy})B.
\end{align}
Because $\sin2\phi=0$ at all four cardinal azimuths, $p_{xy}$ is an exact
geometry null.  If $B$, $A_{xz}$, and $A_{zz}$ are nonzero and absolute
normalization is known, the tensor response rank is four.

\subsection{Common mode and normalization}

The $p_{zz}$ column is proportional to $(1,1,1,1)^{\mathsf T}$.  It is an
independent observable if absolute normalization is externally known.  If a
global normalization is free, the normalization derivative is the same common
mode at the unpolarized point.  After profiling it, the effective tensor rank
of a single ideal ring is three: $(p_{xx}-p_{yy})$, $p_{xz}$, and $p_{yz}$
remain, while $p_{xy}$ is geometric-null and $p_{zz}$ is normalization-null.
Away from the unpolarized point the local null is a state-dependent combination
of scale and tensor common mode, but the physical conclusion is unchanged.

This distinction is the simplest example of why raw response rank and profiled
polarization rank must both be reported.

\section{Multiple polar angles within one station}

For ring $g$ at polar angle $\theta_g$, the common-mode response is
$A_{zz}(\theta_g)/2$.  With a single shared luminosity, two angular settings
separate $p_{zz}$ from luminosity if their normalized response vectors are not
collinear.  In the point-ring limit this requires
\begin{equation}
 A_{zz}(\theta_1)\neq A_{zz}(\theta_2),
 \label{eq:multi-theta-condition}
\end{equation}
and in the real detector it requires unequal acceptance-weighted averages,
\begin{equation}
 \frac{H^{(1)}_{p_{zz}}}{b_1}\neq
 \frac{H^{(2)}_{p_{zz}}}{b_2}.
\end{equation}
If each angular group has its own unconstrained normalization, this recovery is
lost.  Multiple $\theta$ values can also separate vector $A_y$ from tensor
$A_{xz}$ first harmonics when their angular dependences are not proportional.
They cannot recover $p_{xy}$ from perfectly cardinal, azimuth-symmetric arms,
because the $\sin2\phi$ integral remains zero at every polar angle.

The production study evaluates proton and deuteron singles separately and also
evaluates proton--deuteron coincidences.  A union of inclusive singles and
coincidences is deliberately not assigned an independent-Poisson rank because
shared events and shared efficiencies require a consistent exclusive likelihood
or covariance model.

\begin{table}[htbp]
\centering
\caption{Three-hour, unpolarized-point response results for the planned compact
geometry approximation, with one shared profiled global normalization.  Singular
values are derivatives in counts per unit orthonormal tensor coordinate; the
relative/absolute raw-rank thresholds are $10^{-10}/10^{-12}$.}
\scriptsize
\begin{tabularx}{\linewidth}{YYY}
\toprule
Configuration & Raw/profiled rank & Spectrum and dominant nulls\\
\midrule
One $\theta$, four azimuths & $4/3$ analytic & $p_{xy}$ geometry null; $p_{zz}$/scale profiled null\\
Two-$\theta$ proton singles & $4/4$ & $(1.8603,1.8603,1.4074,1.0658,0)\times10^5$; $p_{xy}$ null\\
Deuteron singles, branches summed & $4/3$ & $(5.7183,2.9425,2.9425,0.3166,0)\times10^5$; common mode profiled\\
Deuteron branches, energy resolved & $4/4$ & $(5.0582,2.0818,2.0818,1.7427,0)\times10^5$; $p_{xy}$ null\\
Forward/backward coincidences & $4/4$ & $(1.3879,1.3879,1.0460,0.8852,0)\times10^5$; $p_{xy}$ null\\
Per-angle normalized proton singles & $4/3$ & $p_{zz}$ recovery is lost when each angle has a free scale\\
\bottomrule
\end{tabularx}
\label{tab:scenario-ranks}
\end{table}

\section{Poisson Fisher information and nuisance profiling}

For independent Poisson channels and full parameter vector
$\vect\theta=(\vect q,\vect\eta)$,
\begin{equation}
 F_{ab}=\sum_k\frac{1}{\mu_k}
 \frac{\partial\mu_k}{\partial\theta_a}
 \frac{\partial\mu_k}{\partial\theta_b}.
 \label{eq:fisher}
\end{equation}
Partitioning polarization and nuisance parameters gives
\begin{equation}
 F=
 \begin{pmatrix}
 F_{qq}&F_{q\eta}\\F_{\eta q}&F_{\eta\eta}
 \end{pmatrix}.
\end{equation}
The effective polarization information after nuisance profiling is
\begin{equation}
 F_{\rm eff}=F_{qq}-F_{q\eta}F_{\eta\eta}^{+}F_{\eta q},
 \label{eq:schur}
\end{equation}
where $+$ denotes a tolerance-controlled pseudoinverse.  Gaussian calibration
priors add their precision matrix before the Schur complement.  The output must
include the full Fisher matrix, $F_{\rm eff}$, its eigensystem, covariance
$F_{\rm eff}^{+}$, and identifiable-subspace uncertainties.

At minimum the nuisance model must support a global luminosity scale.  Useful
stress tests add per-angle scales, four relative sector efficiencies, and
analyzing-power modes.  Completely free per-channel efficiencies absorb every
polarization response and eliminate identifiability; calibration priors or
symmetry constraints are therefore part of the measurement definition.

The present scalar code propagates pointwise $T_{20}/T_{22}$ uncertainties in
quadrature rather than fitting analyzing-power nuisance parameters or a table
covariance.  That legacy uncertainty band is not equivalent to
Eq.~\eqref{eq:schur} and must be labeled separately.

\section{Spin transport in the five-dimensional tensor space}

For $R\in SO(3)$,
\begin{equation}
 \vect p_2=R\vect p_1,
 \qquad
 \Tens_2=R\Tens_1R^{\mathsf T}.
 \label{eq:rotation}
\end{equation}
The induced five-dimensional representation in the orthonormal basis is
\begin{equation}
 T^{(2)}_{ab}(R)=\Tr\!\left[E_a^{\mathsf T}RE_bR^{\mathsf T}\right],
 \qquad
 \vect q_2=T^{(2)}(R)\vect q_1.
 \label{eq:rank2-map}
\end{equation}
Because the basis is orthonormal and rotations preserve the Frobenius inner
product, $T^{(2)}(R)$ is orthogonal.  Therefore
\begin{equation}
 \Tr(\Tens_2^2)=\Tr(\Tens_1^2),
 \qquad
 \lVert\vect q_2\rVert_2=\lVert\vect q_1\rVert_2.
\end{equation}
Tests must cover identity, rotations about each Cartesian axis, generic Euler
rotations, composition, and norm preservation.  This layer depends only on
$SO(3)$ and accepts future Thomas--BMT or transfer-map rotations without
embedding a particular magnet lattice.

\section{One station versus two stations}

Let $\vect q_0$ be the tensor at a common reference point and $T_s$ the known
rank--2 transport to station $s$.  The station response to the reference tensor
is $J_sT_s$, so
\begin{equation}
 J_{\rm combined}=
 \begin{pmatrix}
 J_1T_1\\J_2T_2
 \end{pmatrix}.
 \label{eq:combined-response}
\end{equation}
The central geometric statement is
\begin{equation}
 \boxed{
 \ker J_{\rm combined}
 =\ker(J_1T_1)\cap\ker(J_2T_2)}.
 \label{eq:kernel-intersection}
\end{equation}
A second station increases structural rank if and only if this intersection has
smaller dimension than the first station's null space.  Equivalently, at least
one direction invisible to station 1 must have a nonzero projection into the
row space of $J_2T_2$.

Two identical stations with $T_2=T_1$ have aligned null spaces.  Doubling their
exposure improves nonzero Fisher eigenvalues and statistical precision but does
not change rank.  A rotation may still fail to help if it leaves the null
subspace invariant, if the relevant analyzing power vanishes, or if newly
generated response is absorbed by independent station normalizations or
efficiencies.  Thus ``more counts'' and ``more identifiable dimensions'' are
different statements.

With station-specific nuisances, the correct comparison is between profiled
Fisher matrices, not only stacked polarization Jacobians.  Shared luminosity,
common analyzing-power calibration, and independent station efficiencies must
be encoded according to the actual experiment.

\section{Generic rotation sweep and null-space complementarity}

For a relative rotation $R_a(\alpha)$ about $a=x,y,z$, scan
$0\leq\alpha\leq\pi$.  At each point report
\begin{enumerate}
 \item raw and profiled effective rank;
 \item all singular values and the smallest nonzero singular value;
 \item condition number restricted to the identifiable subspace;
 \item Fisher determinant when full rank, otherwise pseudo-determinant;
 \item covariance eigenvalues in the identifiable subspace;
 \item principal angles between transported station null spaces.
\end{enumerate}
If $N_1$ and $N_2$ are orthonormal bases for the two null spaces, the cosines of
their principal angles are the singular values of $N_1^{\mathsf T}N_2$.  A
principal angle of zero identifies a shared null direction.  The dimension of
the exact intersection must agree with Eq.~\eqref{eq:kernel-intersection}.

Extrema in a rotation scan describe mathematical complementarity.  They are not
preferred accelerator settings until a real transport map, acceptance model,
beam phase space, and nuisance constraints are supplied.

\begin{table}[htbp]
\centering
\caption{Generic two-station rotation scan for identical planned proton-channel
geometry at the two stations, a shared global normalization, and a
$\SI{2}{\degree}$ scan grid.  These are mathematical sensitivity rotations, not
an RIBF transfer-map prediction.}
\small
\begin{tabularx}{\linewidth}{YYYY}
\toprule
Rotation family & Unchanged-degeneracy angles & Maximum complementarity & Profiled rank\\
\midrule
$R_x(\alpha)$ & $0^\circ,180^\circ$ & $90^\circ$: $\sigma_{\min}=1.6414\times10^5$ & five for sampled $2^\circ$--$178^\circ$\\
$R_y(\alpha)$ & $0^\circ,180^\circ$ & $90^\circ$: $\sigma_{\min}=1.6414\times10^5$ & five for sampled $2^\circ$--$178^\circ$\\
$R_z(\alpha)$ & $0^\circ,90^\circ,180^\circ$ & $44^\circ/46^\circ$: $\sigma_{\min}=1.0470\times10^5$ & five except listed angles\\
\bottomrule
\end{tabularx}
\label{tab:rotation-summary}
\end{table}

\section{\texorpdfstring{Local analysis around nominal pure $p_{yy}$}{Local analysis around nominal pure pyy}}

An axial tensor with magnitude $P_T$ and unit symmetry axis $\hat n$ is
\begin{equation}
 \Tens(P_T,\hat n)=\frac32P_T\hat n\hat n^{\mathsf T}-\frac12P_TI.
 \label{eq:axial-tensor}
\end{equation}
This normalization is fixed by the repository convention: if $\hat n=\hat y$,
then $p_{yy}=P_T$ and $p_{xx}=p_{zz}=-P_T/2$.

Parameterize small tilts about the vertical direction as
\begin{equation}
 \hat n(\alpha,\beta)=
 \left(\alpha,\sqrt{1-\alpha^2-\beta^2},\beta\right),
\end{equation}
where $\alpha$ tilts the axis toward $+x$ and $\beta$ tilts it toward $+z$.
At $\alpha=\beta=0$,
\begin{align}
 \frac{\partial\Tens}{\partial P_T}
 &=\diag(-1/2,1,-1/2),\\
 \frac{\partial\Tens}{\partial\alpha}
 &=\frac32P_T(\hat x\hat y^{\mathsf T}+\hat y\hat x^{\mathsf T}),\\
 \frac{\partial\Tens}{\partial\beta}
 &=\frac32P_T(\hat y\hat z^{\mathsf T}+\hat z\hat y^{\mathsf T}).
 \label{eq:pyy-derivatives}
\end{align}
Thus the $x$-tilt generates $p_{xy}$ at first order and is exactly invisible to
an ideal cardinal L/R/U/D ring.  The $z$-tilt generates $p_{yz}$ and appears in
the $U-D$ first harmonic when $A_{xz}\neq0$.  Multiple polar angles do not lift
the exact $p_{xy}$ null if every arm remains centered and azimuthally symmetric.
Non-cardinal azimuths, known acceptance asymmetry, or a transported second
station can make that tilt visible; uncontrolled asymmetry can instead mimic it.

The local Fisher analysis in $(P_T,\alpha,\beta)$ must include normalization and
efficiency nuisances.  At exactly $P_T=0$, orientation is undefined and both
tilt derivatives vanish, so an axis ellipse is meaningful only around a
nonzero nominal tensor state.  A second station helps when its transport maps
the upstream $p_{xy}$ perturbation into a tensor harmonic observed by its local
geometry.  Identity transport between identical stations cannot recover the
missing tilt.

Unwanted tensor contamination should be reported by projecting the
five-dimensional covariance onto the three-dimensional tangent space of the
axial model and its two-dimensional orthogonal complement.  This distinguishes
uncertainty in magnitude/axis from non-axial contamination rather than forcing
all deviations into a tilt.

\begin{table}[htbp]
\centering
\caption{Nominal vertical-state local statistical sensitivity at $P_T=0.8$ and
three hours for the planned two-$\theta$ proton channels, profiling one shared
global normalization.  Analyzing-power, background, alignment, and efficiency
systematics are not included.}
\small
\begin{tabularx}{\linewidth}{YcYY}
\toprule
Configuration & $\sigma(P_T)$ & Tilt covariance/ellipse & Contamination eigenmodes\\
\midrule
One station & $0.00378$ & $\alpha$ unidentifiable; $\sigma(\beta)=0.00180$ rad ($0.103^\circ$) & full-tensor $p_{xy}$ null\\
Two stations, $R_y(20^\circ)$ & $0.00267$ & $\sigma(\alpha)=0.00721$ rad, $\sigma(\beta)=0.00180$ rad & no structural tensor null\\
\bottomrule
\end{tabularx}
\label{tab:nominal-pyy}
\end{table}

\section{Detector scenarios and unresolved geometry values}

Several repository artifacts preserve different historical detector values.
The maintained default scalar configuration uses proton laboratory angles
$\SI{53.4}{\degree}$ and $\SI{11.2}{\degree}$, whereas an older paper table
uses $\SI{55.9}{\degree}$ and $\SI{11.3}{\degree}$.  Distances and active sizes
also differ.  The compact common detector configuration currently uses
$\SI{53.4}{\degree}$ at \SI{205}{mm}, $\SI{11.2}{\degree}$ at \SI{190}{mm},
and a deuteron channel at $\SI{20.9}{\degree}$ and \SI{140}{mm}; its detector
response and several mechanical inputs remain provisional.

These values must not be averaged or silently ``corrected.''  Every generated
result must record the exact source configuration, resolved channel manifest,
angle convention, distance reference, active-face dimensions, beam energy,
table checksums, and acceptance algorithm.  Results for the historical paper
geometry and the current compact geometry are different named scenarios.
The after-SRC and pre-SAMURAI compact configurations inherit the same common
physics channel layout at present; their different chamber interfaces alone do
not create a different polarization response matrix.
The CLI selection named \texttt{current} denotes the planned two-angle proton
singles response only.  Deuteron singles and coincidence samples are separate
named studies; inclusive samples are not combined as independent Poisson bins.

\section{Implication for an after-SRC polarimeter}

The relevant conceptual chain is
\begin{equation}
 \begin{gathered}
 \text{PIS}\rightarrow\text{Wien filter}\rightarrow\text{AVF}
 \rightarrow\text{RRC}\rightarrow\text{SRC}\\
 \rightarrow\boxed{\text{after-SRC polarimeter}}
 \rightarrow\text{downstream beam line}\rightarrow\text{SAMURAI}.
 \end{gathered}
\end{equation}
The source and Wien filter prepare the intended spin mode.  The cyclotrons and
extraction system transport that state.  A post-SRC polarimeter independently
anchors the tensor polarization actually delivered by the complete acceleration
chain at final energy.

This argument does not require, and must not be phrased as, an expectation of
large SRC depolarization.  The defensible statement is: good polarization
preservation is expected under proper accelerator operation, while a post-SRC
measurement independently verifies the delivered tensor polarization and
constrains polarization-related systematics.

For the immediate nominal pure-$p_{yy}$ program, one after-SRC station can be
sufficient for the primary task of magnitude and stability monitoring if the
axial-state assumption, angular response, relative efficiencies, and relevant
normalization are controlled.  That statement is narrower than arbitrary
five-component tensor reconstruction.  A second station is scientifically
justified when it addresses an additional quantified question, for example:
\begin{enumerate}
 \item measuring a transport-induced tensor-axis rotation predicted by a known
 beam-line map;
 \item lifting the first-order $p_{xy}$ tilt/contamination null through a
 demonstrably complementary rotation;
 \item separating coherent spin transport from non-axial contamination or
 ensemble depolarization;
 \item providing a downstream polarization anchor when intervening beam-line
 elements introduce a systematic uncertainty relevant at SAMURAI.
\end{enumerate}
The justification must be expressed through the reduction of
$\ker(J_1T_1)\cap\ker(J_2T_2)$ or through a specified improvement in profiled
uncertainty, not merely through doubled event statistics.

\section{Required machine-readable outputs and figures}

For every scenario, the analysis product must include the parameter basis and
conversion matrix, channel metadata, unpolarized vector $\vect b$, response
matrix $H$, full Jacobian, nuisance definitions and priors, column scales,
singular values and vectors, rank thresholds, null vectors in both bases,
Fisher matrices, profiled Fisher matrix, covariance/pseudoinverse, uncertainty
eigenvectors, spin rotation and $T^{(2)}(R)$, and software/configuration
provenance.  JSON is the canonical structured summary and CSV stores matrices
and scans without lossy formatting.

The following figures are referenced only after generation from those products.

\generatedfigure{01_one_station_singular_spectrum.pdf}{Singular-value spectrum for one station.  The caption in the generated figure states the orthonormal basis convention and rank threshold.}{fig:one-station-svd}

\generatedfigure{02_one_station_null_directions.pdf}{Tensor null and near-null directions translated into the physical Cartesian reporting basis.}{fig:null-directions}

\generatedfigure{03_one_theta_vs_multi_theta.pdf}{One-polar-angle versus multiple-polar-angle response, separating structural rank from conditioning.}{fig:multi-theta}

\generatedfigure{04_normalization_profile_comparison.pdf}{Known-normalization and profiled-normalization comparison.  Loss of a common mode must not be labeled a detector-geometry null.}{fig:normalization}

\generatedfigure{05_one_vs_two_station_singular_spectrum.pdf}{One-station and two-station singular spectra for a stated relative spin rotation.}{fig:two-station-svd}

\generatedfigure{06_effective_rank_rotation_scan.pdf}{Effective rank versus generic relative rotation.  This is a sensitivity study, not a beam-line prediction.}{fig:rank-scan}

\generatedfigure{07_smallest_singular_rotation_scan.pdf}{Smallest singular value and threshold versus relative rotation angle.}{fig:smin-scan}

\generatedfigure{08_current_fisher_and_correlation.pdf}{Profiled Fisher/correlation matrix for the planned two-angle proton-singles scenario and nuisance model.}{fig:fisher-correlation}

\generatedfigure{09_nominal_pyy_fisher_ellipses.pdf}{Profiled local covariance for $(P_T,\alpha,\beta)$ around the stated nonzero nominal vertical tensor state.}{fig:pyy-ellipse}

\generatedfigure{10_null_space_intersection_schematic.pdf}{Null-space interpretation of Eq.~\eqref{eq:kernel-intersection}; overlapping null directions remain unobservable after combining stations.}{fig:null-intersection}

\section{Validation requirements}

The tensor layer must test symmetric/traceless round trips, Cartesian/internal
basis round trips, $\Tens'=R\Tens R^{\mathsf T}$, identity rotation, generic
rotation composition, and norm preservation.  The ideal ring test must obtain
the $p_{xy}$ null and the three non-common harmonics from the constructed
matrix, not from hard-coded conclusions.  A normalization test must show the
single-angle profiled-rank loss and conditional recovery with two genuinely
different angular responses.  Two-station tests must show that identical
transport preserves rank, that exposure improves precision without changing
rank, that a selected rotation reduces the null-space intersection, and that a
direct numerical intersection agrees with Eq.~\eqref{eq:kernel-intersection}.

Existing scalar $p_{zz}$ and $p_{yy}$ inference values are regression anchors.
Any change caused by the solid-angle correction, interpolation method, or
geometry resolution requires a focused test and an explicit before/after
record; it must not enter as an incidental side effect of the rank analysis.

\section{Conclusions}

One ideal cardinal L/R/U/D ring supplies three normalization-independent tensor
harmonics: diagonal transverse anisotropy, $p_{xz}$, and $p_{yz}$.  It has an
exact $p_{xy}$ geometry null.  Its $p_{zz}$ common mode is measurable with known
absolute normalization and degenerate with an unknown single-ring scale.
Multiple polar angles can recover $p_{zz}$ when their acceptance-averaged
$A_{zz}$ responses differ under a shared constrained normalization, but cannot
recover $p_{xy}$ while exact cardinal azimuth symmetry remains.

A second station is not automatically a new tensor measurement.  Its value is
determined by the transported null-space intersection in
Eq.~\eqref{eq:kernel-intersection} and by nuisance-aware Fisher information.
For the nominal vertical axial state, magnitude and the tilt generating
$p_{yz}$ are locally accessible to one ideal station, while the orthogonal tilt
generating $p_{xy}$ is first-order invisible.  A known complementary spin
rotation can make that missing direction observable at a second station.

The after-SRC instrument is therefore well motivated as the direct final-energy
anchor for the intended pure-$p_{yy}$ beam.  A second station should be pursued
when a real transport map and a quantified null-space or uncertainty reduction
show that it answers a scientifically distinct downstream question.

\end{document}
